\begin{sourcescontribs}
  The section on segmentally additive metrics is shortly summarized in \textcite[Chapter 12.3.1, pp. 52-57]{foundations2}. A more detailed introduction to metric geometry is \textcite[Chapter 1, pp. 1-64]{busemann}, where curves and segments are the basis for the theory developed afterward. The remark on arc-length is only stated in \textcite[p. 21]{busemann}.\\
  For the section on segmentally additive proximity structures, the main reference was \textcite[Chapter 14.2, pp. 163-174]{foundations2}. In addition to the content provided here, the methods to estimate distances are further explained, and the provided theorems are proved.
  Additional references discussing segmentally additive metrics in the context of proximity structure include \textcite{tversky1970dimensional} and \textcite{foundationsMDS}.
  \cref{ex-minkowski-seg-additive} can be found in \textcite[p. 588]{tversky1970dimensional}.
\end{sourcescontribs}
